% June Week 5 weekly: Mon 2026-06-29 to Sun 2026-07-05.
% This week: leader-facing research summary of v1.8.0 outcome, public
% long-context baseline fact-base, and v2.0.0 decision options.

\section{June Week 5 (Mon 2026-06-29 to Sun 2026-07-05)}

\begin{frame}{Reporting Period}
  \begin{center}
    \textbf{Plan 08 / Paper B --- June Week 5} \\[0.3em]
    \small Mon 2026-06-29 to Sun 2026-07-05 \\
    \scriptsize Leader update: research outcome, long-context facts, v2.0.0 options, and resource ask.
  \end{center}
\end{frame}

\begin{frame}{Scientific Question}
  \small
  \begin{block}{Main question}
    When a model faces long evidence, what should it keep, what should it ignore,
    and how can it know whether its reduced context is still reliable?
  \end{block}
  \begin{block}{Why this is scientific, not only engineering}
    ``Long context'' is not one capability. It decomposes into retrieval,
    local extraction, distributed semantic reasoning, distractor robustness, and
    architecture-dependent memory access.
  \end{block}
  \begin{itemize}
    \item The core object of study is the \textbf{interaction between task type, context length, and reduction method}.
    \item The scientific goal is to map this interaction, then design a method
      that adapts instead of assuming one fixed rule.
  \end{itemize}
\end{frame}

\begin{frame}{Background: Qwen3.5 / 3.6 Hybrid GDN Stack}
  \scriptsize
  \begin{center}
  \begin{tikzpicture}[scale=0.82, transform shape,
    node distance=0.35cm,
    box/.style={draw, rounded corners, align=center, minimum height=0.55cm, font=\scriptsize},
    gdn/.style={box, fill=blue!10, minimum width=1.35cm},
    attn/.style={box, fill=orange!18, minimum width=1.35cm},
    state/.style={box, fill=green!12, minimum width=2.2cm},
    kv/.style={box, fill=red!10, minimum width=2.2cm},
    note/.style={align=left, font=\scriptsize},
    arr/.style={-{Latex[length=2mm]}, thick}
  ]
    \node[box, fill=gray!10, minimum width=1.6cm] (tok) {input\\tokens};
    \node[gdn, right=0.65cm of tok] (g1) {GDN\\layer};
    \node[gdn, right=0.25cm of g1] (g2) {GDN\\layer};
    \node[gdn, right=0.25cm of g2] (g3) {GDN\\layer};
    \node[attn, right=0.25cm of g3] (a1) {gated\\attention};
    \node[gdn, right=0.25cm of a1] (g4) {GDN\\layers};
    \node[attn, right=0.25cm of g4] (a2) {gated\\attention};
    \node[box, fill=gray!10, right=0.65cm of a2, minimum width=1.4cm] (out) {reader\\hidden};

    \draw[arr] (tok) -- (g1);
    \draw[arr] (g1) -- (g2);
    \draw[arr] (g2) -- (g3);
    \draw[arr] (g3) -- (a1);
    \draw[arr] (a1) -- (g4);
    \draw[arr] (g4) -- (a2);
    \draw[arr] (a2) -- (out);

    \node[state, below=1.0cm of g2] (s) {bounded recurrent state\\history compressed into state};
    \node[kv, below=1.0cm of a1] (kv1) {KV cache only for\\attention layers};
    \node[kv, below=1.0cm of a2] (kv2) {KV cache only for\\attention layers};

    \draw[arr, blue!60] (g1.south) -- (s.north west);
    \draw[arr, blue!60] (g2.south) -- (s.north);
    \draw[arr, blue!60] (g3.south) -- (s.north east);
    \draw[arr, orange!70] (a1.south) -- (kv1.north);
    \draw[arr, orange!70] (a2.south) -- (kv2.north);

    \node[note, below=1.05cm of s, text width=0.86\textwidth] {
      \textbf{Prefill dependency:} most layers scan tokens into a GDN recurrent state; periodic attention layers
      create ordinary K/V for exact lookup. \textbf{KV-cache methods apply only to the attention subset, not to GDN layers.}
      Qwen3.6 is reported as KV on 16/64 layers, so the KV lever is partial.
    };
  \end{tikzpicture}
  \end{center}
\end{frame}

\begin{frame}{Background: MiniMax-Text-01 / M1 Hybrid Lightning Stack}
  \scriptsize
  \begin{center}
  \begin{tikzpicture}[scale=0.82, transform shape,
    node distance=0.28cm,
    box/.style={draw, rounded corners, align=center, minimum height=0.55cm, font=\scriptsize},
    lin/.style={box, fill=blue!10, minimum width=1.02cm},
    soft/.style={box, fill=orange!18, minimum width=1.15cm},
    moe/.style={box, fill=purple!10, minimum width=1.1cm},
    state/.style={box, fill=green!12, minimum width=2.3cm},
    kv/.style={box, fill=red!10, minimum width=2.3cm},
    note/.style={align=left, font=\scriptsize},
    arr/.style={-{Latex[length=2mm]}, thick}
  ]
    \node[box, fill=gray!10, minimum width=1.15cm] (tok) {input};
    \node[lin, right=0.55cm of tok] (l1) {LA};
    \node[lin, right=0.12cm of l1] (l2) {LA};
    \node[lin, right=0.12cm of l2] (l3) {LA};
    \node[lin, right=0.12cm of l3] (l4) {LA};
    \node[lin, right=0.12cm of l4] (l5) {LA};
    \node[lin, right=0.12cm of l5] (l6) {LA};
    \node[lin, right=0.12cm of l6] (l7) {LA};
    \node[soft, right=0.2cm of l7] (sa) {softmax\\attn};
    \node[moe, right=0.32cm of sa] (moe) {MoE\\FFN};
    \node[box, fill=gray!10, right=0.55cm of moe, minimum width=1.2cm] (out) {next\\block};

    \draw[arr] (tok) -- (l1);
    \draw[arr] (l1) -- (l2);
    \draw[arr] (l2) -- (l3);
    \draw[arr] (l3) -- (l4);
    \draw[arr] (l4) -- (l5);
    \draw[arr] (l5) -- (l6);
    \draw[arr] (l6) -- (l7);
    \draw[arr] (l7) -- (sa);
    \draw[arr] (sa) -- (moe);
    \draw[arr] (moe) -- (out);

    \node[state, below=1.05cm of l4] (s) {Lightning / linear-attention state\\cheap long-range scan};
    \node[kv, below=1.05cm of sa] (kv) {KV cache only for\\periodic softmax block};

    \draw[arr, blue!60] (l2.south) -- (s.north west);
    \draw[arr, blue!60] (l4.south) -- (s.north);
    \draw[arr, blue!60] (l6.south) -- (s.north east);
    \draw[arr, orange!70] (sa.south) -- (kv.north);

    \node[note, below=1.05cm of s, text width=0.86\textwidth] {
      \textbf{Prefill dependency:} the published MiniMax-Text-01/M1 pattern is roughly
      \textbf{7 Lightning/linear-attention blocks + 1 softmax-attention block}, with MoE capacity.
      Lightning blocks carry cheap long-range state; softmax blocks periodically restore exact token interactions.
      \textbf{KV eviction applies only to the softmax blocks.}
    };
  \end{tikzpicture}
  \end{center}
\end{frame}

\begin{frame}{Background: MiniMax-M3 Sparse-Attention Stack}
  \scriptsize
  \begin{center}
  \begin{tikzpicture}[scale=0.82, transform shape,
    node distance=0.35cm,
    box/.style={draw, rounded corners, align=center, minimum height=0.55cm, font=\scriptsize},
    dense/.style={box, fill=orange!18, minimum width=1.45cm},
    sparse/.style={box, fill=cyan!12, minimum width=1.55cm},
    idx/.style={box, fill=yellow!18, minimum width=2.2cm},
    kv/.style={box, fill=red!10, minimum width=2.2cm},
    note/.style={align=left, font=\scriptsize},
    arr/.style={-{Latex[length=2mm]}, thick}
  ]
    \node[box, fill=gray!10, minimum width=1.25cm] (tok) {input};
    \node[dense, right=0.7cm of tok] (d1) {full\\attention};
    \node[sparse, right=0.35cm of d1] (sp1) {MSA\\sparse};
    \node[dense, right=0.35cm of sp1] (d2) {full\\attention};
    \node[sparse, right=0.35cm of d2] (sp2) {MSA\\sparse};
    \node[box, fill=gray!10, right=0.7cm of sp2, minimum width=1.25cm] (out) {output};

    \draw[arr] (tok) -- (d1);
    \draw[arr] (d1) -- (sp1);
    \draw[arr] (sp1) -- (d2);
    \draw[arr] (d2) -- (sp2);
    \draw[arr] (sp2) -- (out);

    \node[kv, below=1.0cm of d1] (kv1) {K/V blocks exist\\for attention layers};
    \node[idx, below=1.0cm of sp1] (idx1) {indexer scores key blocks\\top blocks + local blocks};
    \node[kv, below=1.0cm of sp2] (kv2) {sparse access to\\selected KV blocks};

    \draw[arr, orange!70] (d1.south) -- (kv1.north);
    \draw[arr, cyan!70] (sp1.south) -- (idx1.north);
    \draw[arr, cyan!70] (idx1.east) -- (kv2.west);
    \draw[arr, cyan!70] (sp2.south) -- (kv2.north);

    \node[note, below=1.05cm of idx1, text width=0.86\textwidth] {
      \textbf{Prefill dependency:} M3 is not recurrent/linear. It is a sparse-attention Transformer:
      sparse layers use an indexer to choose key blocks, then exact attention runs on those selected blocks.
      \textbf{KV exists, but access is already learned and sparse}; KV eviction methods are possible only as a comparison to MSA, not as the native mechanism.
    };
  \end{tikzpicture}
  \end{center}
\end{frame}

\begin{frame}{Experimental Facts We Established}
  \small
  \begin{itemize}
    \item \textbf{Fact 1: fixed baselines are regime-specific.} Window, RAG,
      prompt compression, and KV eviction each work in some regimes and fail in others.
    \item \textbf{Fact 2: length changes method ranking.} Attention-based KV
      eviction works around shorter contexts but can collapse at 16k+.
    \item \textbf{Fact 3: task type changes method ranking.} Retrieval,
      extractive QA, global reasoning, and literary MC prefer different context
      reduction behavior.
    \item \textbf{Fact 4: more context is not always better.} Full context can
      hurt when the context introduces distractors or exceeds the model's usable reasoning ability.
    \item \textbf{Fact 5: model architecture matters.} Linear-attention / GDN-style
      bases do not expose the same KV-cache lever as dense-attention models.
  \end{itemize}
\end{frame}

\begin{frame}{Current Paper Intro: Story We Can Tell}
  \small
  \begin{block}{Intro logic}
    Modern LLMs can accept longer contexts, but long-context use is not solved:
    full context is costly and sometimes harmful; simple reduction methods are brittle;
    and different tasks require different evidence-preservation strategies.
  \end{block}
  \begin{itemize}
    \item Start with the field assumption: longer windows and context compression should help.
    \item Show the contradiction: the same method can win on retrieval and lose on semantic QA; full context can even hurt.
    \item Therefore the missing piece is not ``one better compressor'' but a
      mechanism that can \textbf{adapt to the evidence regime}.
    \item This motivates learned memory plus a reliability decision.
  \end{itemize}
\end{frame}

\begin{frame}{What v1.8.0 Already Gives Us}
  \small
  \begin{itemize}
    \item A working learned-memory path on high-headroom tasks after fixing
      termination supervision and joint answer signal.
    \item A gate/fallback formulation that is useful conceptually: use memory when
      reliable, fall back when unsafe.
    \item Ablation evidence that the carrier is not fragile: many knobs move the
      compressed path by only a few points.
    \item A clearer negative result: the current gate story must become held-out
      or conformal, not in-sample threshold picking.
  \end{itemize}
  \begin{block}{Takeaway}
    v1.8.0 is the prototype and evidence base. v2.0.0 should sharpen the claim
    into a cleaner reliability theory and a better necessity regime.
  \end{block}
\end{frame}

\begin{frame}{Next Paper Initialization Direction}
  \small
  \begin{block}{Working title direction}
    Adaptive Evidence Memory for Long-Context Reasoning
  \end{block}
  \begin{itemize}
    \item \textbf{Paper role of the fact-base:} establish that fixed black-box
      context reduction is not enough.
    \item \textbf{Paper role of the method:} introduce a memory that preserves
      evidence and a reliability signal that decides when to trust it.
    \item \textbf{Core novelty target:} self-verified memory + calibrated routing,
      not just another compressor.
    \item \textbf{Evaluation target:} public long-context regimes where window,
      RAG, prompt compression, and KV eviction fail for different reasons.
  \end{itemize}
\end{frame}

\begin{frame}{v2.0.0 Ideas}
  \small
  \begin{enumerate}
    \item \textbf{Self-verified evidence memory:} train the memory to preserve evidence;
      use that preservation/reconstruction signal as the reliability check.
    \item \textbf{Reliability-gated fallback:} use extracted knowledge when safe;
      otherwise fall back to raw context, retrieval, or a higher-cost path.
    \item \textbf{Hybrid evidence memory:} combine a raw recent window with compressed global memory.
    \item \textbf{Adaptive budget:} easy contexts get fewer memory tokens; dense or ambiguous contexts get more.
  \end{enumerate}
\end{frame}

\begin{frame}{Next Experiments To Initialize}
  \small
  \begin{itemize}
    \item \textbf{M1: Hybrid Evidence Memory} --- reader sees raw recent window + compressed global memory.
    \item \textbf{M2: Conformal router} --- held-out reliability calibration instead of in-sample threshold picking.
    \item \textbf{M3: Long-context necessity wave} --- semantic-spread tasks plus retrieval honest-negative.
    \item \textbf{M4: Transfer matrix} --- train on one domain, test on another; measure whether memory is domain-invariant.
    \item \textbf{M5: Figures} --- turn the fact-base into length, task, and method-failure figures.
  \end{itemize}
\end{frame}

\begin{frame}{What Leaders Need to Decide}
  \small
  \begin{description}
    \item[Option A: Fact-base paper] Publish the long-context baseline map first; method comes later.
    \item[Option B: Method paper] Build v2.0.0 now around self-verified evidence memory + conformal routing.
    \item[Option C: Two-stage plan] Fact-base as Section 1 / motivation, then v2.0.0 as the method contribution.
  \end{description}
  \begin{block}{Recommendation}
    Option C: use the fact-base to define the gap, then spend compute only on the
    minimum v2.0.0 experiments needed to show the gap can be closed.
  \end{block}
\end{frame}

\begin{frame}{Recommended Next Step}
  \small
  \begin{itemize}
    \item Freeze the public long-context fact-base and produce the figures.
    \item Implement Hybrid Evidence Memory: raw window + compressed memory.
    \item Add conformal reliability routing and report risk--coverage / cost--coverage.
    \item Run the 5x5 transfer matrix to test cross-domain memory generalization.
    \item Decide whether v2.0.0 is a full method paper or a fact-base + position paper.
  \end{itemize}
  \begin{block}{Ask}
    Approve GPU budget for the v2.0.0 public benchmark experiments, especially
    the hybrid-memory and transfer-matrix runs.
  \end{block}
\end{frame}
